\documentclass[11pt]{article}

% ---------------------------------------------------------------------------
% Packages (arXiv-friendly: all in TeX Live)
% ---------------------------------------------------------------------------
\usepackage[utf8]{inputenc}
\usepackage[T1]{fontenc}
\usepackage{lmodern}
\usepackage[margin=1in]{geometry}
\usepackage{amsmath,amssymb}
\usepackage{booktabs}
\usepackage{array}
\usepackage{graphicx}
\usepackage[table]{xcolor}
\usepackage[round,authoryear]{natbib}
\usepackage{microtype}
\usepackage[hidelinks]{hyperref}
\usepackage{url}

% ---------------------------------------------------------------------------
% Drafting macros
% ---------------------------------------------------------------------------
% \todo{...}: a result or decision that is still pending.
\DeclareRobustCommand{\todo}[1]{\textcolor{red}{\textbf{[TODO}\if\relax\detokenize{#1}\relax\else: #1\fi\textbf{]}}}
% \citeneeded{...}: a claim that needs a reference we have not yet verified.
\DeclareRobustCommand{\citeneeded}[1]{\textcolor{orange}{[citation needed: #1]}}
% Short placeholder for table cells.
\newcommand{\tc}{\textcolor{red}{\textbf{?}}}

\newcommand{\WIS}{\mathrm{WIS}}
\newcommand{\Gemb}{G_{\mathrm{emb}}}
\newcommand{\Gid}{G_{\mathrm{id}}}
\newcommand{\Gex}{G_{\mathrm{excess}}}
\newcommand{\Cstar}{C^{\ast}}

\title{Place or Placeholder?\\[0.3em]
\large When satellite embeddings genuinely improve dengue forecasts}
\author{Author names TBD}
\date{Draft of \today}

\begin{document}
\maketitle

% ===========================================================================
\begin{abstract}
Recent studies report that learned place embeddings improve sub-national dengue forecasts.
Those evaluations used a single embedding snapshot, did not compare the embeddings with
location-only features, and tested only on districts seen in training. Because an embedding
that differs between districts can act as a district identifier, such gains could reflect
memorised location rather than information about the environment. We test this for the
AlphaEarth Foundations satellite embedding (64 dimensions, annual) in weekly dengue
forecasting for 6{,}990 second-level administrative units in Brazil, Colombia, Peru and Taiwan
(OpenDengue v1.3.1). We fit a feature ladder from a seasonal-naive baseline to case history,
climate, conventional remote sensing and population, and finally the AlphaEarth embedding of the
year before the target. Every embedding is aligned in time, and we compare it against five
location-only controls: target encoding, spherical harmonics of coordinates, shuffled embeddings,
random spatial fields and a static 2023 snapshot. Evaluation covers temporal, spatial-block,
held-out-state, held-out-unit and held-out-country splits. Forecasts are scored with the weighted
interval score (WIS) on $\log(1+\text{cases})$. The primary estimand is the excess gain
$\Gex = \WIS(L3{+}\Cstar) - \WIS(L4{+}\Cstar)$, the improvement from embeddings beyond the best
location control. In Brazil under a temporal split, LightGBM with case history reaches a log-WIS
of 0.141--0.199 at 1--4 weeks ahead, and adding ERA5-Land climate lowers it by only 0.3--0.8\%
(not significant after Holm correction). The best location-only control, chosen on a 2019
tuning origin, is a leakage-free target encoding. \todo{Headline $\Gemb$, $\Gid$ and $\Gex$ with
CIs for Brazil S1 and the spatial splits; cross-country result; one-sentence conclusion.}
\end{abstract}

% ===========================================================================
\section{Introduction}
\label{sec:intro}

Dengue forecasting at the district scale supports vector control and clinical preparedness,
but week-to-week case counts are dominated by the district's own history and by season.
Two recent studies report that adding place embeddings from Google's Population Dynamics
Foundation Model \citep[PDFM;][]{agarwal2024pdfm} improves probabilistic dengue forecasts:
for 5{,}570 Brazilian municipalities \citep{brittain2026foundation}, and for about 2{,}450 Mexican
municipalities in one case study of a broader paper \citep{muslim2026planetary}. A related study
reports gains from the AlphaEarth satellite embedding \citep{brown2025alphaearth} for several
climate-sensitive health outcomes \citep{nazir2026alphaearth}.

These results raise a measurement question. An embedding is a fixed vector per district (or per
district-year). A flexible forecaster can use any such vector as an identifier: it can learn
that a given vector belongs to a district with high baseline incidence, without using anything
the vector says about land cover, urban form or climate. Under a temporal split, where every test
district also appears in training, this identifier role is enough to improve scores. We call
this the \emph{location-identifier confound}. It is testable in three ways. First, replace the
embedding with features that carry location and nothing else, and see whether they reproduce the
gain. Second, evaluate on districts, regions or countries the model has never seen, where an
identifier cannot help. Third, align the embedding in time, so that a forecast uses only
imagery from before the target year.

The published dengue evaluations do not yet do these things. \citet{brittain2026foundation} use
a single October 2023 PDFM snapshot throughout a 2020--2025 backtest, with temporal hold-outs only.
We found no location-only control in their text. The Mexico dengue analysis of
\citet{muslim2026planetary} also uses a static October 2023 snapshot as its primary analysis (with
an October 2025 sensitivity check). It has no location-only control and no held-out region,
although the same paper uses shuffled-embedding controls and held-out jurisdictions in other case
studies. The reported dengue gain there is small: a mean WIS change of $-0.0051$ at one month,
which is not significant at three or six months, and the median municipality was slightly worse.
\citet{nazir2026alphaearth} use AlphaEarth for malaria as a static per-location fingerprint. Their
respiratory-infection analysis is year-aligned, and neither analysis includes a location-only
control. The closest precedent is a city-scale study in India \citep{pradhan2026benchmarks}. Adding
embeddings to forecasters left error metrics largely unchanged. In a small-area mapping task, the
embedding advantage over open map layers fell to 0.001 once the layers were given each cell's
coordinates. The author attributes the null results to the small number of independent spatial
units at city scale.

\paragraph{Contributions.}
We provide a multi-country test of whether satellite embeddings add forecast skill beyond
location. Specifically:
\begin{enumerate}
  \item A pre-specified evaluation protocol for weekly district dengue forecasts across four
  countries and 6{,}990 units. It uses a feature ladder (L0--L4) with every level trained on the same rows, so
  the embedding comparison is not confounded by training-window length.
  \item Five location-only controls (C1--C5) that replace the embedding block. They separate
  ``any district-specific vector'', ``smooth map of position'' and ``future imagery''
  explanations of an embedding gain.
  \item Six train/test splits (S0--S5), from a deliberately leaky random split to leave-one-country-out.
  They measure how any gain behaves when the identifier route is closed.
  \item A difference estimand, $\Gex$, that is stable when the embedding gain is near zero. The
  ``ID share'' ratio is reported only when the embedding gain is clearly non-zero.
  \item A documented, validated data pipeline: zero-filling of OpenDengue with external checks,
  boundary joins at $\geq 99.99\%$ of cases, and exact, verified population-weighted AlphaEarth
  district means.
\end{enumerate}
The analysis is designed so that either outcome is informative. If gains vanish under the
controls, published improvements should be read as location effects. If gains survive,
the splits and strata (RQ5 below) indicate where imagery adds information.

\paragraph{Research questions.}
RQ1: do embeddings improve forecasts under the temporal split used by earlier work (S1, L4 vs L3)?
RQ2: how much of that gain does a location-only control reproduce?
RQ3: does any gain survive on unseen units, regions and countries (S2--S5)?
RQ4: does a time-aligned embedding differ from a static 2023 snapshot (C5)?
RQ5: where does imagery help most (urban growth, case-history length, population)?

% ===========================================================================
\section{Related work}
\label{sec:related}

\paragraph{Place embeddings for disease forecasting.}
PDFM \citep{agarwal2024pdfm} provides 330-dimensional embeddings derived from search, maps and
activity data. It covers 17 countries, including Brazil but not Colombia, Peru or Taiwan.
\citet{brittain2026foundation} add PDFM as a post-hoc residual adjustment to five forecasting
families (SARIMA, XGBoost, TCN, N-BEATS and zero-shot TimesFM). They forecast monthly
municipal dengue in Brazil at 1--12 months and report improved WIS in 19 of 20 model-horizon
combinations. \citet{muslim2026planetary} report five PDFM case studies, one of which is
Mexican municipal dengue with TimesFM 2.0 plus a PDFM residual model, WIS, Diebold--Mariano
tests and a block bootstrap. Our model stack (gradient boosting, TimesFM, WIS, DM) overlaps with
theirs on purpose, so that our S1 results can be compared with theirs. The difference lies in the
controls, the time alignment and the transfer splits.

\paragraph{Satellite embeddings and location.}
AlphaEarth Foundations \citep{brown2025alphaearth} produces a 64-dimensional, unit-length embedding
per 10\,m pixel per year from 2017. Its authors themselves compare against controls for purely
geographic information. Location encoders built from coordinates are a natural control for any
embedding \citep{russwurm2024location,klemmer2025satclip}. For
population estimation, PDFM gains are uneven across space and scale \citep{zhang2026population}.
\citet{pradhan2026benchmarks} report that, in one city, satellite embeddings were not
distinguishable from open map layers once coordinates were added.

\paragraph{Spatial validation.}
Random cross-validation inflates skill for spatially structured data. Blocked and spatial designs
are the standard remedy \citep{roberts2017cv,ploton2020spatial}, and predictions outside the
training feature space need explicit assessment \citep{meyer2021aoa,meyer2022global}. Our S2--S5
splits apply this to forecasting, and S0 is the leaky random contrast.

\paragraph{Scoring epidemic forecasts.}
We use the weighted interval score \citep{bracher2021wis}, computed on the log scale as recommended
for epidemic counts \citep{bosse2023transformed}. Forecasts are compared with the Diebold--Mariano test
\citep{diebold1995comparing} with the small-sample correction of \citet{harvey1997testing}.
Zero-shot TimesFM \citep{das2024timesfm} is now a common baseline. We are not aware of a
peer-reviewed WIS evaluation of TimesFM 2.5 for infectious-disease counts.
The closest Brazilian benchmark, the 2024 Dengue Forecasting Sprint, scored state-level forecasts
with the log score, CRPS and interval score \citep{araujo2025sprint}.

% ===========================================================================
\section{Data}
\label{sec:data}

\subsection{Dengue cases: OpenDengue v1.3.1}
We use the Temporal extract of OpenDengue v1.3.1 \citep{clarke2024opendengue,opendengue131},
restricted to rows with weekly time resolution and second-level (Admin2) spatial resolution. An
independent check confirmed that the three release archives match the GitHub blob hashes of release commit 99b6a3e
and the figshare v5 MD5 sums, and that they parse without error
(Appendix~\ref{app:audit}). All weekly rows are 7-day spans starting on Sunday, and no (unit, week)
pair is duplicated. A unit is identified by its IBGE municipal code in Brazil and by its
Admin1|Admin2 name elsewhere. FAO GAUL codes are not unique: Taiwan's 287 units share two codes.
Table~\ref{tab:panel} summarises the panel.

\begin{table}[t]
\centering\small
\caption{Weekly district panel after zero-filling (OpenDengue v1.3.1, Temporal extract,
Admin2 $\times$ Week). Units are those that ever reported a case. ``Analysis years'' are the
target years that the protocol uses.}
\label{tab:panel}
\begin{tabular}{lrrlrll}
\toprule
Country & Units & Weeks & Weekly coverage & \% zero-filled & Analysis years & Case definition \\
\midrule
Brazil   & 5{,}517 & 574   & 2013-12-29 to 2024-12-22 & 40.1 & 2014--2024 & Probable \\
Colombia & 1{,}071 & 835   & 2006-12-31 to 2022-12-25 & 77.7 & $\le$2022 & Total \\
Peru     & 115     & 1{,}252 & 2000-01-02 to 2023-12-24 & 78.5 & $\le$2023 & Probable and confirmed \\
Taiwan   & 287     & 1{,}409 & 1997-12-28 to 2024-12-22 & 98.1 & $\le$2024 (no 2021 cases) & Local confirmed$^{a}$ \\
\bottomrule
\end{tabular}

\smallskip
\footnotesize $^{a}$OpenDengue labels Taiwan ``Total''. Our cross-check shows that the series equals locally
acquired confirmed cases by onset week and residence township (Section~\ref{sec:zerofill}).
\end{table}

The unit counts are units that \emph{ever reported a case}, not administrative totals. Peru's
115 units are provinces (of 196), and Taiwan's 287 are townships (of about 368) that ever had a
locally acquired case. Peru is therefore used only for the temporal split (S1) and the
cross-country split (S5). Taiwan is extremely sparse: 3.7\% of its unit-week grid has a row, and most
cases fall in the 2014, 2015 and 2023 outbreaks. We treat it as a stress test, report it per
outbreak and at county level, and do not use it for headline claims.

\subsection{Zero-filling and its validation}
\label{sec:zerofill}
OpenDengue omits most zero-case weeks, and inconsistently. Brazil lists explicit zeros for active
municipalities in 2014--2021 but none from 2022. Colombia, Peru and Taiwan list case-positive
weeks only. Without a fix, the reporting format would change inside the test period. We
therefore freeze the following rule. The unit universe is every Admin2 unit that appears in a
country's weekly file. Every week inside the country's weekly coverage window that has no row
is set to zero and flagged \texttt{zero\_filled}.

We validated the rule against two external sources (Appendix~\ref{app:zerofill}):
\begin{itemize}
  \item \textbf{Brazil against InfoDengue.} \citep{codeco2018infodengue,infodengue_api} We sampled 40
  municipalities in 2017, 2019, 2022 and 2023, including municipality-years absent from
  OpenDengue, municipalities present in only some weeks, and capitals. OpenDengue matches InfoDengue's
  \emph{probable} case count (\texttt{casprov}) closely. The ratio of sums is 0.94--1.02, the median
  weekly difference is 0 or $-1$, and the weekly correlation is $\geq 0.97$ in every group. It does
  not match notified cases. In absent weeks, probable counts were zero in every
  sampled week of 2017 and 2019. In 2022--2023, 1.1--2.1\% of absent weeks had 1--8 probable cases,
  which is consistent with later revisions of the surveillance data. 2024 was not checked.
  \item \textbf{Taiwan against the Taiwan CDC case file.} \citep{taiwancdc_dengue_daily} OpenDengue's Taiwan series equals locally acquired confirmed cases by onset week and residence
  township. Weekly national totals match exactly in 83.7\% of weeks with $r=1.0000$. All 7{,}673 present unit-weeks match exactly, and
  no unit-week absent from OpenDengue has a local case. The missing year 2021 had no local cases.
\end{itemize}
We also compared Admin2 sums with OpenDengue's national weekly series. Colombia agrees closely.
In Brazil, Admin2 (probable) is a smooth 0.3--0.7 of the national series (total) from 2016 onward,
which is a case-definition offset rather than missing weeks. A week is flagged as a reporting gap (\texttt{flag\_gap})
when its Admin2/national ratio falls below one quarter of that country-year's ratio. Only one
week is flagged (Peru, 2018-08-12), and it is excluded from scoring.

\subsection{Boundaries}
Brazilian units are joined by IBGE code to the IBGE 2022 municipal mesh \citep{ibge2022malha}.
Colombia, Peru and Taiwan are joined by name, within Admin1, to geoBoundaries gbOpen ADM2
polygons \citep{geoboundaries_gbopen}. Matching runs in order through manual fixes, exact,
normalised and fuzzy matches. For Taiwan it also uses a township crosswalk inferred from the CDC case file,
because the official township file was not reachable. The join covers 6{,}987 of 6{,}990 units
and at least 99.99\% of 2014+ cases in every country (Appendix~\ref{app:boundaries}). The
unmatched units are Bel\'en de Bajir\'a (Colombia, 40 cases, municipality status annulled) and
two non-spatial ``other'' buckets in Taiwan. An independent verifier re-downloaded every source,
reviewed every fuzzy, manual and sampled crosswalk match, and found no suspect matches.

\subsection{Covariates}
\label{sec:covariates}
All gridded inputs are fetched from Earth Engine \citep{gorelick2017gee} with
\texttt{computePixels} and reduced locally to polygons with exact coverage fractions
\citep{baston_exactextract}. Population weights come from GHSL GHS-POP \citep{ghsl_pop}, with
negative no-data values masked to zero.
\begin{itemize}
  \item \textbf{Climate (L2).} ERA5-Land daily aggregates \citep{era5land_ee} are summarised to the panel's
  Sunday-start weeks. The variables are mean, minimum and maximum 2\,m temperature, precipitation, and
  relative humidity from temperature and dewpoint. They are population-weighted, with the GHSL epoch
  nearest the year. Local reduction agrees with EE \texttt{reduceRegions} to $r>0.99999999$ at about
  1/20 of the compute. All countries are complete over their panel windows. Three units with no ERA5-Land land cell use the
  nearest land cell.
  \item \textbf{Remote sensing (L3).} Monthly MODIS NDVI \citep{mod13q1_ee} (QA-masked, read
  at about 1\,km; $r=0.998$ against native 250\,m) and day and night land-surface temperature
  \citep{mod11a2_ee}. Each composite is assigned to the month of its last day, so month $m$ contains
  no information from after month $m$. Gaps are filled from the same month of the previous year only.
  Annual VIIRS night lights \citep{viirs_ee} and GHSL built-up share \citep{ghsl_built}.
  \item \textbf{Population (L3).} GHSL population interpolated log-linearly between epochs.
  It exceeds the 2022 Brazilian census by 4--10\%, so we use it as a relative scale, not a count.
  WorldPop \citep{worldpop_r2025a} is reserved for a sensitivity check.
\end{itemize}

\subsection{AlphaEarth district embeddings}
\label{sec:aef}
We use the AlphaEarth Foundations Satellite Embedding V1 annual collection
\citep{aef_dataset,brown2025alphaearth}: 64 bands (A00--A63), 10\,m, with unit-length pixel
vectors. Earth Engine holds the years 2017--2025. For each unit and image year we compute two means:
\begin{itemize}
\item a population-weighted mean at 100\,m on the native GHS-POP Mollweide grid, with GHSL
interpolated linearly between the 2015, 2020 and 2025 epochs;
\item an unweighted (area) mean at 300\,m.
\end{itemize}
A pilot on 100 Brazilian municipalities showed that 100\,m is effectively exact for
population-weighted means (mean cosine 0.99994 against 10\,m), whereas naive 300\,m resampling loses
accuracy (mean cosine 0.993).

Extraction fetches embedding blocks with \texttt{computePixels} in float32, only over populated
blocks, and accumulates weighted sums locally. This was exact (cosine 0.9999998 against
\texttt{reduceRegions}) at about 40\% of the cost of \texttt{reduceRegions}, and was the cheapest
of the eight estimators we piloted (Appendix~\ref{app:aef}).

District means are not unit-length. Population-weighted norms have a median of about 0.92
(range 0.74--0.98), and the shrinkage varies by unit. We therefore renormalise each mean to unit
length, and keep the raw norm for a sensitivity analysis.

\paragraph{Verification.}
An independent script recomputed 33 Brazilian unit-years with its own Earth Engine code. The sample
spans small, medium, large, coastal and Amazon units in all seven years. It found
a minimum cosine of 0.9999999987 (weighted) and 0.9999999969 (unweighted), and population sums
within 0.02\%. Completeness is 5{,}517/5{,}517 units in every year. The one
exception is Fernando de Noronha, which has no valid AlphaEarth pixels in 2017--2021 upstream, so its
embedding features are missing for target years 2018--2022. Spatially adjacent units have a median
cosine of 0.973 against 0.710 for random pairs, and consecutive years have a median cosine of 0.98.
Re-extracting one state-year with a different degree of parallelism gave bitwise-identical sums.
\todo{Verification summaries for Colombia, Peru and Taiwan.}

% ===========================================================================
\section{Methods}
\label{sec:methods}

\subsection{Forecasting task}
The target is weekly reported cases per unit, modelled as $y=\log(1+\text{cases})$. Forecasts are
issued every week for horizons $h=1,\dots,4$ weeks, as nine quantiles (0.1, 0.2, \dots, 0.9). We use
rolling-origin annual refits. For test year $Y\in\{2020,\dots\}$ (Brazil to 2024, Colombia 2022,
Peru 2023, Taiwan 2020 and 2022--2024), models are trained on target weeks in years 2018 to $Y-1$
and forecast every target week of $Y$. Lagged inputs may reach back to 2013, or to the country's first weekly year.
\textbf{Every ladder level is trained on the same 2018+ rows}, because AlphaEarth starts in 2017
and the embedding of year $y-1$ is first available for 2018 targets. Hyperparameters are tuned once
on the 2019 origin (train 2018, validate 2019) with the same budget per level, then frozen. A
secondary analysis trains L0--L3 on the long window (2013 to $Y-1$).

\subsection{Feature ladder}
Table~\ref{tab:ladder} lists the levels. L0--L2 contain no location-specific features, and the code enforces this
by assertion. L3 deliberately adds unit-level covariates: population and built-up share
vary between units and could themselves act as weak identifiers. All weekly inputs are read
relative to the forecast origin and checked to precede it. Monthly inputs use only months that
ended before the origin week. Yearly inputs (night lights, built-up share, population, AlphaEarth)
for a target in epidemiological year $y$ come from year $y-1$.

\begin{table}[t]
\centering\small
\caption{Feature ladder. $p$ is the number of columns in Brazil.}
\label{tab:ladder}
\begin{tabular}{llr}
\toprule
Level & Adds & $p$ \\
\midrule
L0 & Same week one year earlier (seasonal naive) & 1 \\
L1 & Case lags 0--11 from the origin; target week $-52j$ ($\pm 1$ week), $j=1,2,3$; week-of-year sine/cosine & 23 \\
L2 & + ERA5-Land lags 0--11 (temperature, precipitation, RH) and 4/8/12-week rolling means of 5 variables & 74 \\
L3 & + NDVI, day and night LST (last 3 complete months); night lights, built-up share, log population ($y-1$) & 86 \\
L4 & + AlphaEarth population-weighted mean of year $y-1$, renormalised (64) & 150 \\
L5 & + PDFM static snapshot (Brazil only; optional) \todo{access and inclusion} & -- \\
\bottomrule
\end{tabular}
\end{table}

\subsection{Models}
\begin{itemize}
  \item \textbf{Seasonal naive (L0).} The quantiles of the training-set errors
  $\log(1+y_T)-\log(1+y_{T-52})$ are added to last year's value.
  \item \textbf{Negative binomial (L1).} An NB2 GLM on raw counts with the L1 features and an intercept, fitted by L-BFGS.
  Predictive quantiles come from the NB distribution. It is not tuned.
  \item \textbf{LightGBM quantile regression} \citep{ke2017lightgbm} \textbf{(main model).}
  One booster per quantile and horizon, on $\log(1+\text{cases})$. Fixed settings: learning rate 0.05,
  400 rounds, feature fraction 0.9, bagging 0.8, $\lambda_2=1$, deterministic mode. The grid is
  \texttt{num\_leaves} $\in\{31,127\}$ $\times$ \texttt{min\_data\_in\_leaf} $\in\{100,1000\}$,
  selected by mean pinball loss at $\tau\in\{0.1,0.5,0.9\}$, $h\in\{1,4\}$ on the 2019 origin.
  This selected 127/100 for L1 and L2, and 31/100 for L3, where the grid was almost flat. \todo{L4 tuned parameters.}
  Quantile crossings are removed by sorting.
  \item \textbf{TimesFM 2.5, zero-shot} \citep{das2024timesfm,timesfm25}. Input is each unit's
  $\log(1+\text{cases})$ history up to the origin (up to 573 weeks; maximum context 576), with the continuous quantile head.
  It natively outputs the mean and the deciles 0.1--0.9, which fixes our primary quantile set.
  A hybrid in which LightGBM corrects TimesFM residuals using the ladder features is planned
  \todo{TimesFM + LightGBM residual hybrid}.
\end{itemize}

\subsection{Location-only controls}
\label{sec:controls}
Each control \emph{replaces} the AlphaEarth block of L4, giving a model ``L3+C$k$''. Controls are
rebuilt for every split, fold, origin and seed.
\begin{description}
  \item[C1: target encoding.] Each unit's mean $\log(1+\text{cases})$. Under S0/S1 the mean
  is taken over panel years \emph{before the first training target year} (Brazil 2014--2017), so it is causal and
  identical for training and test rows. An earlier version took the mean over the training targets,
  which leaked each training row's target into its own feature (Appendix~\ref{app:refinements}).
  For unseen units (S2--S4), C1 is an inverse-distance mean over the five nearest training units in
  the same country. Training units receive a leave-one-out value, so that training and test encodings match.
  C1 is not defined across countries (S5).
  \item[C2: coordinates.] Latitude, longitude and orthonormal real spherical harmonics of degree
  1--8 (80 terms, no constant).
  \item[C3: shuffled embeddings.] Real AlphaEarth vectors permuted across units within each
  country-year. Under S2--S4, a held-out unit receives a random training unit's vector. 10 seeds.
  \item[C4: random smooth fields.] 64 Gaussian random fields per country with exponential covariance,
  whose range is fitted to the empirical semivariogram of the AlphaEarth district means. 10 seeds.
  \item[C5: static snapshot.] The AlphaEarth 2023 embedding for every target year. This mirrors the
  static-snapshot design of earlier work and tests RQ4.
\end{description}
Units are located by a polygon representative point, which avoids offshore centroids. The best control
$\Cstar$ is the one with the lowest mean log-WIS on the 2019 origin, chosen before any test-year
comparison of L3 and L4.

\subsection{Splits}
\begin{description}
  \item[S0 (random; leaky).] A hash of (unit, week) sends 20\% of 2018..$Y$ rows to test. Reported
  only as an upper-bound contrast (Brazil, Colombia, Taiwan).
  \item[S1 (temporal).] All units in training and test. This is the setting of earlier work.
  \item[S2 (held-out state).] Leave one Admin1 out (Brazil 27 folds, Colombia 33).
  \item[S3 (spatial blocks with buffer).] $k$-means blocks on unit locations (Brazil 10, Colombia 6). Training
  units within distance $D$ of the test block are dropped. $D$ is the first distance at which the
  correlogram of de-seasonalised residuals, estimated on 2014--2019, falls below 0.1: 225\,km for Brazil and 130\,km for Colombia
  (Appendix~\ref{app:refinements}).
  \item[S4 (held-out units).] Five folds of units, stratified by Admin1.
  \item[S5 (held-out country).] Train on three countries and test on the fourth, Taiwan first.
  Cross-country comparisons are additionally scored on log incidence per 100{,}000 \todo{per-country scaling
  rule}.
\end{description}
Held-out units keep their own lagged cases as inputs. ``Cold start'' means that the model never
trained on the unit's rows, not that the unit has no history.

\subsection{Scoring}
The primary score is the weighted interval score \citep{bracher2021wis} on $\log(1+\text{cases})$
\citep{bosse2023transformed}, using the median and the 20/40/60/80\% central intervals formed by
the nine deciles:
\[
\WIS = \frac{1}{K+1/2}\Bigl(\tfrac12|y-m| + \sum_{k=1}^{K}\tfrac{\alpha_k}{2}\,
\mathrm{IS}_{\alpha_k}\Bigr),\quad
\mathrm{IS}_\alpha = (u-l) + \tfrac{2}{\alpha}(l-y)_+ + \tfrac{2}{\alpha}(y-u)_+ .
\]
Scores are averaged over unit-weeks within each country $\times$ split $\times$ horizon, and every
model is scored on an identical set of rows. Secondary metrics are WIS on the case scale and on
incidence, 23-quantile WIS for models that produce it, coverage of the 40\% and 80\% intervals,
and median absolute error. We also report the share of units improved and the median per-unit difference,
and stratify by observed versus zero-filled targets. Outbreak detection is defined with outbreak weeks
as cases $\geq\max(\text{unit 90th percentile over 2013--2019},\,5)$, scored as sensitivity at a 10\% false-alarm
rate \todo{outbreak analysis}. A reporting-lag sensitivity check drops the last two weeks of history
\todo{reporting-lag check}.

\subsection{Estimands}
With positive values meaning lower (better) WIS:
\begin{align*}
\Gemb &= \WIS(L3) - \WIS(L4) &&\text{embedding gain},\\
\Gid  &= \WIS(L3) - \WIS(L3{+}\Cstar) &&\text{location-control gain},\\
\Gex  &= \WIS(L3{+}\Cstar) - \WIS(L4{+}\Cstar) &&\text{embedding gain beyond location (primary)}.
\end{align*}
The ratio $\mathrm{ID\ share}=\Gid/\Gemb$ is reported only if the 95\% interval of $\Gemb$
excludes zero \emph{and} $\Gemb\geq 1\%$ of $\WIS(L3)$. Otherwise it is reported as ``not estimable''.
RQ4 compares L4 with L3+C5, and RQ5 stratifies $\Gemb$ and $\Gex$ by GHSL urban growth, case-history
length and population.

\subsection{Inference}
For each comparison and horizon, we average the per-unit WIS difference over units within each
forecast week to obtain a weekly series $d_t$. We then apply the Diebold--Mariano test
\citep{diebold1995comparing} with the Harvey--Leybourne--Newbold correction
\citep{harvey1997testing}. The long-run variance is the equal-weighted cosine (EWC) estimator with
$B=\lfloor 0.4\,T^{2/3}\rfloor$ and fixed-$b$ $t_B$ critical values \citep{lazarus2018har}. We chose EWC over Newey--West because, in a
simulation with AR(1) loss differences ($\phi=0.5$, $T=104$), Newey--West had 10.8\% size at nominal
5\% and EWC 6.9\%. Confidence intervals come from a moving-block bootstrap
\citep{kunsch1989jackknife} \citeneeded{circular block bootstrap, Politis and Romano 1992,
in Exploring the Limits of Bootstrap; chapter pages not verified}
that resamples Admin1 clusters crossed with circular 8-week blocks, using the same resamples for all
estimands. $p$-values are adjusted by Holm's method \citep{holm1979simple} across the pre-specified
primary comparisons \todo{state the Holm family as frozen}.

\subsection{Pre-specification}
The splits, estimands, control-selection rule, quantile set, scale and inference were fixed in a design
document before the primary test-year comparisons \todo{freeze date (planned 15 January 2027) and
archive link}. Test-year scores of the baselines (L0--L2, TimesFM) were inspected during pipeline
development and are reported as such \todo{confirm the status of L3/L4 test-year results at freeze}.

% ===========================================================================
\section{Results}
\label{sec:results}

\subsection{Baselines and the lower ladder (Brazil, S1)}
Table~\ref{tab:br-s1} gives the Brazil temporal results. All models are scored on the same
1{,}439{,}937 rows per horizon (5{,}517 units, 2020--2024). LightGBM with case history (L1) is the
best of the baselines at every horizon and in every test year. It beats L0, NB and zero-shot TimesFM in all
Diebold--Mariano tests, with Holm-adjusted $p\leq 0.0015$ over 16 tests ($T=261$ weeks). Its mean
log-WIS is lower than L0's in 5{,}516 of 5{,}517 units. On the case scale, zero-shot TimesFM has a
lower WIS than LightGBM (1.85 vs.\ 2.32 at $h=1$), and the NB GLM is unstable: a few
epidemic-peak forecasts, e.g.\ S\~ao Paulo in 2024, dominate its case-scale mean.

Adding ERA5-Land climate (L2) lowers log-WIS by 0.0004 ($h=1$) to 0.0016 ($h=4$), i.e.\ 0.3--0.8\%.
The raw DM $p$-values are 0.03--0.08, and none survives Holm correction ($p_{\text{Holm}}=0.109$
at every horizon). L2 is better than L1 in every test year at $h=4$. Its whole gain comes from observed
weeks, not zero-filled ones. Once twelve weeks of case lags and three seasons of seasonal lags are
present, climate adds little at 1--4 weeks.

\begin{table}[t]
\centering\small
\caption{Brazil, S1: mean WIS on $\log(1+\text{cases})$ (nine deciles), pooled over test years 2020--2024.
Lower is better. \textcolor{red}{Red} cells are pending.}
\label{tab:br-s1}
\begin{tabular}{llrrrr}
\toprule
Model & Level & $h=1$ & $h=2$ & $h=3$ & $h=4$ \\
\midrule
Seasonal naive & L0 & 0.4645 & 0.4645 & 0.4645 & 0.4645 \\
Negative binomial & L1 & 0.1894 & 0.2102 & 0.2279 & 0.2454 \\
TimesFM 2.5 zero-shot & -- & 0.1594 & 0.1829 & 0.2051 & 0.2252 \\
LightGBM & L1 & 0.1414 & 0.1628 & 0.1813 & 0.1992 \\
LightGBM & L2 & 0.1410 & 0.1621 & 0.1801 & 0.1975 \\
LightGBM & L3 & \tc & \tc & \tc & \tc \\
LightGBM & L4 & \tc & \tc & \tc & \tc \\
LightGBM & L3+$\Cstar$ (C1) & \tc & \tc & \tc & \tc \\
LightGBM & L4+$\Cstar$ (C1) & \tc & \tc & \tc & \tc \\
LightGBM & L3+C5 (static AE 2023) & \tc & \tc & \tc & \tc \\
TimesFM + LightGBM residual & L4 & \tc & \tc & \tc & \tc \\
\bottomrule
\end{tabular}
\end{table}

\subsection{Choice of the location control}
Table~\ref{tab:cstar} shows the 2019 selection origin for Brazil. With a single training year,
C2--C5 all score slightly \emph{worse} than plain L3. The leakage-free target encoding C1 is the only
control that improves on L3 ($-0.00062$ log-WIS pooled), and it is lowest at $h=2$--4. Its margin
over the next control (C3, shuffled embeddings) is 0.00088, about nine times the seed standard deviation of C3.
Hence $\Cstar=\text{C1}$ for Brazil S1. The C4 range fitted to the Brazilian embeddings
(8{,}483\,km) is weakly identified, because the semivariogram rises almost linearly with no sill.
\todo{$\Cstar$ for Colombia, Peru, Taiwan; for Taiwan, selection is dominated by zero targets.}

\begin{table}[t]
\centering\small
\caption{Selection of $\Cstar$ on the 2019 origin (Brazil, LightGBM with frozen L3 parameters;
train 2018, validate 2019; 1{,}147{,}536 rows per run). Mean log-WIS pooled over $h=1$--4. C3 and C4: mean over
10 seeds.}
\label{tab:cstar}
\begin{tabular}{lrrrrrr}
\toprule
Run & pooled & seed SD & $h=1$ & $h=4$ & cov 80 & vs.\ L3 \\
\midrule
L3+C1 (pre-training-years encoding) & \textbf{0.17462} & -- & 0.14281 & 0.20423 & 0.888 & $-0.00062$ \\
L3 (reference) & 0.17524 & -- & 0.14283 & 0.20572 & 0.889 & 0 \\
L3+C3 (shuffled AE) & 0.17550 & 0.00010 & 0.14298 & 0.20621 & 0.887 & $+0.00026$ \\
L3+C5 (AE 2023 static) & 0.17599 & -- & 0.14321 & 0.20681 & 0.884 & $+0.00075$ \\
L3+C2 (spherical harmonics) & 0.17647 & -- & 0.14339 & 0.20767 & 0.885 & $+0.00123$ \\
L3+C4 (random fields) & 0.17649 & 0.00017 & 0.14336 & 0.20781 & 0.884 & $+0.00126$ \\
\midrule
\textit{L3+C1, training-target encoding (leaky; superseded)} & \textit{0.26314} & -- & \textit{0.22454} & \textit{0.29789} & \textit{0.815} & \textit{$+0.08791$} \\
\bottomrule
\end{tabular}
\end{table}

\subsection{Embedding gain and excess gain}
Table~\ref{tab:estimands} is the main result table. \todo{Fill from the L3, L4, L3+$\Cstar$, L4+$\Cstar$
runs; describe RQ1 (S1 $\Gemb$), RQ2 ($\Gid$, ID share where estimable) and RQ3 (S2--S5).}

\begin{table}[t]
\centering\small
\setlength{\tabcolsep}{4pt}
\caption{Estimands by country and split, LightGBM, pooled over horizons \todo{or per horizon}.
Gains are in log-WIS units (positive = embedding or control helps), with 95\% block-bootstrap
intervals. $p$: Holm-adjusted DM $p$-value for $\Gex$. ID share is shown only when estimable.}
\label{tab:estimands}
\begin{tabular}{llcccccc}
\toprule
Country & Split & $\Cstar$ & $\WIS(L3)$ & $\Gemb$ [95\% CI] & $\Gid$ [95\% CI] & $\Gex$ [95\% CI] & $p$ / ID share \\
\midrule
Brazil   & S0 & \tc & \tc & \tc & \tc & \tc & \tc \\
         & S1 & C1 & \tc & \tc & \tc & \tc & \tc \\
         & S2 & \tc & \tc & \tc & \tc & \tc & \tc \\
         & S3 & \tc & \tc & \tc & \tc & \tc & \tc \\
         & S4 & \tc & \tc & \tc & \tc & \tc & \tc \\
         & S5 & \tc & \tc & \tc & \tc & \tc & \tc \\
\midrule
Colombia & S0 & \tc & \tc & \tc & \tc & \tc & \tc \\
         & S1 & \tc & \tc & \tc & \tc & \tc & \tc \\
         & S2 & \tc & \tc & \tc & \tc & \tc & \tc \\
         & S3 & \tc & \tc & \tc & \tc & \tc & \tc \\
         & S4 & \tc & \tc & \tc & \tc & \tc & \tc \\
         & S5 & \tc & \tc & \tc & \tc & \tc & \tc \\
\midrule
Peru     & S1 & \tc & \tc & \tc & \tc & \tc & \tc \\
         & S5 & \tc & \tc & \tc & \tc & \tc & \tc \\
\midrule
Taiwan   & S0 & \tc & \tc & \tc & \tc & \tc & \tc \\
         & S1 & \tc & \tc & \tc & \tc & \tc & \tc \\
         & S5 & \tc & \tc & \tc & \tc & \tc & \tc \\
\bottomrule
\end{tabular}

\smallskip
\footnotesize S5 rows are scored on the held-out country. C1 is undefined under S5, so $\Cstar$ is chosen among C2--C5.
\end{table}

\subsection{Time alignment (RQ4)}
\todo{L4 vs.\ L3+C5 by test year. Note that for 2024 targets C5 and L4 use the same 2023 embedding.}

\subsection{Calibration}
Table~\ref{tab:coverage} reports interval coverage. Coverage over all targets looks adequate only
because 73.8\% of targets are zero. On non-zero targets every model under-covers: the 80\% interval
contains the truth in 59--76\% of cases. LightGBM's 80\% coverage on non-zero targets is 0.64--0.69 in every year,
and climate does not change this.

\begin{table}[t]
\centering\small
\caption{Empirical coverage of the central 40\% and 80\% intervals, Brazil S1, 2020--2024.
Non-zero targets are 26.2\% of rows.}
\label{tab:coverage}
\begin{tabular}{lcccc}
\toprule
& \multicolumn{2}{c}{All targets} & \multicolumn{2}{c}{Cases $>0$} \\
\cmidrule(lr){2-3}\cmidrule(lr){4-5}
Model & 40\% & 80\% & 40\% & 80\% \\
\midrule
Seasonal naive (L0) & 0.643 & 0.773 & 0.043 & 0.400 \\
NB (L1) & 0.829 & 0.937 & 0.355 & 0.762 \\
LightGBM L1 & 0.781 & 0.907 & 0.312 & 0.677 \\
LightGBM L2 & 0.774 & 0.904 & 0.306 & 0.667 \\
TimesFM 2.5 ZS & 0.617 & 0.810 & 0.267 & 0.593 \\
LightGBM L3 / L4 & \tc & \tc & \tc & \tc \\
LightGBM L3+$\Cstar$ / L4+$\Cstar$ & \tc & \tc & \tc & \tc \\
\bottomrule
\end{tabular}
\end{table}

\subsection{Strata (RQ5)}
Zero-filled targets make up 26--32\% of Brazilian test rows in 2020--2021 and 61--74\% in 2022--2024.
Table~\ref{tab:strata} therefore separates observed from zero-filled targets alongside the RQ5 strata.
On zero-filled rows the NB GLM is best (0.026), and LightGBM's error there rises from about 0.002
(2020--21) to 0.044--0.066 (2022--24). Year-to-year comparisons of pooled means partly reflect the
zero-fill share.

\begin{table}[t]
\centering\small
\caption{Mean log-WIS by stratum, Brazil S1 (all horizons).}
\label{tab:strata}
\begin{tabular}{lcccc}
\toprule
Stratum & L1 LightGBM & L2 LightGBM & $\Gemb$ & $\Gex$ \\
\midrule
Observed targets & 0.3181 & 0.3159 & \tc & \tc \\
Zero-filled targets & 0.0407 & 0.0408 & \tc & \tc \\
GHSL urban growth, tercile 1/2/3 & \tc & \tc & \tc & \tc \\
Case-history length, tercile 1/2/3 & \tc & \tc & \tc & \tc \\
Population, tercile 1/2/3 & \tc & \tc & \tc & \tc \\
Outbreak weeks & \tc & \tc & \tc & \tc \\
\bottomrule
\end{tabular}
\end{table}

\subsection{Other countries, monthly arm and sensitivity analyses}
\todo{Colombia/Peru/Taiwan S1 ladders; Taiwan per outbreak (2014 Kaohsiung, 2015 Tainan/Kaohsiung,
2023 Tainan) and at county level; monthly replication arm (Brazil, S1, L3/L4/L5, 1 and 3 months);
long-window L0--L3; unweighted and non-renormalised embeddings; LST v2 (QC-masked); WorldPop weights.}

% ===========================================================================
\section{Discussion}
\label{sec:discussion}

\todo{Interpretation once $\Gex$ is available. Write separately for each outcome: (a) $\Gex\approx 0$ with
$\Gemb>0$: an identifier effect; (b) $\Gex>0$ in S1 only: the information is unit-specific but does not
transfer; (c) $\Gex>0$ under S2--S5: environmental information.}

Several findings already constrain the interpretation. First, under a temporal split the
history-only LightGBM baseline is strong. It improves on the seasonal-naive forecast in essentially
every Brazilian municipality, and climate adds less than 1\% at 1--4 weeks. Any embedding gain must
therefore be measured against a baseline that leaves little between-unit level information unexplained.
Second, a correctly causal target encoding was the only location control that improved on L3
at the selection origin. A naive in-sample target encoding instead made forecasts much worse (+0.088
log-WIS) through target leakage. Location controls must therefore be built with the same care as
the features they are compared against. Third, annual embeddings vary only between units and
slowly across years, never within a season. Any gain they provide must come from differences in
level or in seasonal amplitude between units, which is the information a location identifier can also
supply. This is why the controls and the spatial splits are central rather than supplementary.

Relative to the PDFM studies, our design differs in embedding source (satellite imagery only),
temporal resolution (weekly, 1--4 weeks), and the explicit time alignment of the embedding.
Our S1 results are therefore not a replication of \citet{brittain2026foundation} or
\citet{muslim2026planetary}. The monthly arm and the Brazil-only L5 are intended to give a closer
comparison \todo{}.

% ===========================================================================
\section{Limitations}
\label{sec:limitations}

\begin{itemize}
  \item \textbf{Taiwan is thin.} Most Taiwanese units report no local case in most years. At the 2019
  tuning origin only 95 (train) and 56 (validation) non-zero $h=1$ targets exist among about 14{,}900 rows.
  Tuning and $\Cstar$ selection there are nearly uninformative, and Taiwan results are mainly the 2023
  outbreak. Peru has 115 provinces, so its spatial resolution is coarse.
  \item \textbf{Surveillance definitions and revisions.} Case definitions differ by country (probable,
  total, local confirmed). OpenDengue is a snapshot: in Brazil 2022--2023, 1--2\% of zero-filled weeks
  had a few probable cases in later InfoDengue data, and 2024 was not cross-checked. The 2022 change
  in Brazil's zero reporting means that 61--74\% of 2022--2024 test targets are zero-filled.
  \item \textbf{Reporting lag.} Models are trained and scored on finalised data, not on what was known
  on each forecast date. A two-week drop sensitivity is planned but is not a full real-time evaluation.
  \item \textbf{Land-surface temperature artefacts.} MODIS LST without a QC mask lets cloud-contaminated
  cold pixels through in cloudy, mountainous units: Colombian night LST reaches $-42\,^\circ$C. No MOD11A2 QC rule
  fixed this cleanly: a strict mask shifted all countries by $+0.2$ to $+1.5\,^\circ$C and created new
  outliers. We use the unmasked series with a causal screen (outside Brazil), and the masked series
  as a sensitivity check.
  \item \textbf{Population.} GHSL interpolated between epochs exceeds the 2022 Brazilian census by
  4--10\%, and its 2025 epoch is a projection. Population enters as a relative weight and scale.
  Interpolating between epochs also lets a later epoch inform earlier years' weights.
  \item \textbf{Yearly-feature alignment.} Yearly inputs for target year $y$ come from year $y-1$.
  For the few January targets whose origin falls late in year $y-1$, the annual product overlaps the
  weeks after the origin. Operationally, an annual embedding is published some time after its year ends,
  so L4 as specified is not exactly a real-time forecast. The AlphaEarth model itself was trained after
  the study period. Time alignment of its inputs does not rule out information from that training
  \todo{check what the AlphaEarth documentation states about training years}.
  \item \textbf{Boundary vintages.} Polygons are from 2017--2022 sources (Taiwan from OpenStreetMap 2017,
  because the official township file was unavailable), and the Taiwan township crosswalk was inferred
  from case counts. It reproduces every OpenDengue total but is not an official code list.
  \item \textbf{Scope of controls.} The controls test specific alternative explanations. A non-significant
  $\Gex$ does not show that embeddings carry no environmental information, only that it does not
  improve these forecasts beyond location at 1--4 weeks.
\end{itemize}

% ===========================================================================
\section*{Data and code availability}
OpenDengue v1.3.1 is available from its GitHub release and figshare \citep{opendengue131} under
CC BY-SA. All covariates are public Earth Engine collections. The AlphaEarth Foundations Satellite
Embedding dataset is produced by Google and Google DeepMind (CC-BY 4.0). Derived district-week tables
will be released under CC BY-SA 4.0, as required by OpenDengue's and the Taiwan boundaries' share-alike
terms. Code (data pipeline, controls, splits, runner, scoring and statistics, with unit tests) will be
released under an open licence \todo{repository URL, licence file, Zenodo DOI}. Exact package versions are
recorded with every run.

% ===========================================================================
\bibliographystyle{plainnat}
\bibliography{refs}

% ===========================================================================
\appendix
\clearpage

\section{OpenDengue audit}
\label{app:audit}
\textbf{Integrity.} The National, Spatial and Temporal zips match the GitHub blob hashes at commit 99b6a3e
(3 November 2026), the figshare v5 MD5 sums, and their extracted CSVs. All files parse fully,
with no invalid dates, negative or non-integer counts, or duplicate rows. Data end on 2025-04-30.
Quirks handled: literal \texttt{"NA"} strings, \texttt{confirmed}/\texttt{Confirmed}, and four India
rows with an embedded line break.

\textbf{Resolution.} The Temporal extract is a strict superset of the Spatial extract for weekly
Admin2 data. It adds 306 Peruvian unit-weeks, and no value conflicts.

\begin{table}[h]
\centering\small
\caption{Weekly Admin2 audit by country.}
\begin{tabular}{lrrrrl}
\toprule
Country & Units & GAUL codes & Weekly rows & Rows with 0 & Grid completeness \\
\midrule
Brazil & 5{,}517 & 5{,}456 & 1{,}895{,}538 & 63\% & 59.9\% (2013: 1 week only) \\
Colombia & 1{,}071 & 1{,}012 & 199{,}458 & 0\% & 22.3\% \\
Peru & 115 & 113 & 30{,}991 & 0\% & 21.5\% \\
Taiwan & 287 & 2 & 7{,}673 & 0\% & 3.7\% (2021 absent) \\
\bottomrule
\end{tabular}
\end{table}

\textbf{Fallback countries.} Mexico has Admin1 data only, and Thailand has no weekly data, so neither can
replace a district-level country.

\section{Zero-fill cross-check}
\label{app:zerofill}

\begin{table}[h]
\centering\small
\caption{Brazil: OpenDengue versus InfoDengue probable cases (\texttt{casprov}), 40 sampled
municipalities.}
\begin{tabular}{lrrrr|rrr}
\toprule
& \multicolumn{4}{c|}{Present weeks} & \multicolumn{3}{c}{Absent weeks} \\
Year & $n$ & median diff. & $\Sigma$OD/$\Sigma$casprov & exact (capitals) & $n$ & casprov$>0$ & max \\
\midrule
2017 & 1{,}218 & 0 & 1.008 & 58\% & 902 & 0 & 0 \\
2019 & 1{,}507 & 0 & 1.003 & 52\% & 573 & 0 & 0 \\
2022 & 795 & 0 & 1.020 & 22\% & 1{,}285 & 14 (1.1\%) & 4 \\
2023 & 849 & $-1$ & 0.939 & 8\% & 1{,}231 & 26 (2.1\%) & 8 \\
\bottomrule
\end{tabular}
\end{table}

Taiwan: for 1998--2024, OpenDengue sums to 101{,}892 cases against 102{,}302 locally acquired CDC cases by onset week.
The 410-case difference is early local cases with no recorded residence county (217 in 2002). Yearly totals match
exactly in the outbreak years (2014: 15{,}506; 2015: 43{,}467; 2023: 26{,}462). The CDC-to-OpenDengue township
crosswalk was inferred by maximising shared cases within county, and it reproduces 100\% of OpenDengue cases.

\section{Boundary joins}
\label{app:boundaries}

\begin{table}[h]
\centering\small
\caption{Boundary join rates (cases are weekly Admin2 cases from 2014).}
\begin{tabular}{llrrrr}
\toprule
Country & Source & Units & Matched & Cases & Case \% \\
\midrule
Brazil & IBGE 2022 (code join) & 5{,}517 & 5{,}517 & 16{,}708{,}965 & 100.000 \\
Colombia & geoBoundaries (DANE MGN 2020) & 1{,}071 & 1{,}070 & 687{,}419 & 99.994 \\
Peru & geoBoundaries (IGN/OCHA) & 115 & 115 & 578{,}638 & 100.000 \\
Taiwan & geoBoundaries (OSM 2017) & 287 & 285 & 86{,}360 & 100.000 \\
\bottomrule
\end{tabular}
\end{table}

Colombia: 1{,}005 exact, 30 normalised, 28 fuzzy and 7 manual matches. Taiwan: 271 via the CDC crosswalk,
6 manual, 4 exact, 3 normalised and 1 fuzzy. No polygon is matched to more than one unit, and all geometries are
valid. Three unit centroids lie offshore or outside their polygon (Vit\'oria, Putumayo, Qijin), so
representative points are used for distances.

\section{AlphaEarth extraction and verification}
\label{app:aef}

\begin{table}[h]
\centering\small
\caption{Estimator pilot (Brazil, 2022). Cosine of the population-weighted mean against the 100\,m
reference (mean / min), and projected cost for 7 years $\times$ 4 countries.}
\begin{tabular}{lccr}
\toprule
Estimator & cosine (mean / min) & EECU-s per km$^2$ & EECU-h \\
\midrule
100\,m \texttt{reduceRegions} (reference) & 1 / 1 & 0.025 & 516 \\
Population-proportional points, $N=250$ & 0.99947 / 0.9972 & -- & 3{,}900 \\
300\,m pyramid & 0.9921 / 0.9516 & 0.0039 & 81 \\
200\,m pyramid & 0.9981 / 0.9888 & 0.0075 & 154 \\
\textbf{\texttt{computePixels} 100\,m float32 + local zonal} & \textbf{0.9999998 / 0.9999989} & \textbf{0.0100} & \textbf{207} \\
\texttt{computePixels} 100\,m int8 & 0.999996 / 0.999994 & 0.0130 & 268 \\
\bottomrule
\end{tabular}
\end{table}

Production fetched only populated 128$\times$128 blocks at 100\,m, at 0.0119 EECU-s per fetched km$^2$, and
used a 300\,m pass for unweighted means. The cost was about 21 EECU-h per Brazil-year. Production verification
against fresh \texttt{reduceRegions} gave cosine 1.0000000 at eight decimals and population sums within 0.008\%. The independent
verification is summarised in Section~\ref{sec:aef}. Year-to-year cosine is lowest for 2017$\to$2018
along the north-eastern coast (minimum 0.74). These units recompute exactly, so the shift is in the
upstream product.

\section{Compute and determinism}
\label{app:compute}
\begin{itemize}
  \item Earth Engine: climate for all countries took about 3.6 EECU-h, remote sensing about 8.9 EECU-h (plus 8.65 EECU-h for
  the QC-masked LST variant), and AlphaEarth about 270 EECU-h projected for 2017--2024 across countries
  \todo{final ledger totals}.
  \item LightGBM: one Brazil L3 chunk (1.44\,M training rows, nine quantile boosters) takes about 8 minutes at 16 threads.
  The full protocol is about 41{,}000 core-hours at 4 threads, 72\% of it in the C3/C4 seed replicates
  \todo{final seed policy for spatial splits}.
  \item Predictions were bit-identical across repeated runs and across 2, 4, 8 and 16 threads. Compared runs
  nevertheless use identical thread counts, LightGBM version (4.7.0), deterministic mode and seeds.
  \item TimesFM 2.5 zero-shot: 1.52\,M series for Brazil in about 22 minutes on one 48\,GB GPU, with 8.1\,GB peak memory.
\end{itemize}

\section{Implementation refinements to the protocol}
\label{app:refinements}
\begin{itemize}
  \item \textbf{DM variance.} We use the EWC default. Null rejection rates at nominal 5\% over 1{,}000 AR(1) replicates, $h=1$:
  $\phi=0.5$, $T=104$: EWC 6.9\%, Newey--West 10.8\%, Bartlett lag $h-1$ 27.5\%. $\phi=0$, $T=104$: EWC 5.8\%.
  $\phi=0.5$, $T=260$: EWC 5.7\%.
  \item \textbf{S3 buffer.} On protocol residuals (unit and week-of-year effects removed), the Brazilian correlogram plateaus at
  about 0.11 out to 2{,}000\,km, which reflects national co-movement. $D$ is therefore the maximum of the 0.1 crossings
  with and without removing the national weekly mean: 225\,km for Brazil, 130\,km for Colombia.
  \item \textbf{C1 leakage.} Encoding each unit by its mean over the training targets placed each training row's
  target inside its own feature. At the 2019 origin this gave log-WIS 0.263 against 0.175 for L3. The pre-training-years
  encoding removes the leak, and tests confirm that C1 is unchanged when every row from the first training year
  onward is overwritten.
  \item \textbf{C4 range.} Brazil: 8{,}483\,km, from the 2017--2018 embeddings (2017 alone 11{,}398\,km, 2018 alone
  6{,}105\,km). Colombia 705\,km, Peru 148\,km, Taiwan 27\,km.
  \item \textbf{LST screen (Colombia, Peru, Taiwan).} Unit-months more than 10\,$^\circ$C below the unit's
  calendar-month median for years $\leq 2017$ are set to missing. This flags 0.57\% of Colombian night values and
  0.06--0.18\% elsewhere.
  \item \textbf{Missing embeddings.} Fernando de Noronha (2017--2021) has missing \texttt{ae\_*} values. LightGBM routes
  missing values natively, and C5 (2023) is defined for it.
\end{itemize}

\end{document}
